\chapter{Visão geral e objetivos}

\textbf{O HEP Conference Hub }é um portal destinado a localizar, organizar, acompanhar e auditar eventos científicos relacionados à Física de Altas Energias (HEP), calorimetria, detectores, aceleradores, computação científica, cosmologia e astropartículas. O sistema foi projetado para reduzir a dispersão de informações entre páginas oficiais, plataformas Indico, sites de institutos, agendas regionais e páginas de colaboração.

O portal procura responder, em uma única interface, às seguintes perguntas:

\begin{itemize}[leftmargin=*,itemsep=0.25em]
\item Quais eventos estão próximos, em andamento ou concluídos?
\item A submissão de resumos está aberta, encerrada ou ainda não anunciada?
\item A inscrição está aberta e quais são os prazos antecipado, regular ou presencial?
\item Qual é a página exata para submissão, inscrição, programa, cronograma, contribuições ou proceedings?
\item Quando cada informação foi verificada, qual foi a fonte e qual o grau de confiança?
\item Quais eventos pertencem à América Latina e quais são apenas candidatos ainda não validados?
\end{itemize}

\section{Público-alvo}

\begin{itemize}[leftmargin=*,itemsep=0.25em]
\item Pesquisadores, docentes e estudantes de pós-graduação em HEP/HECAP;
\item Membros de colaborações experimentais e grupos de instrumentação;
\item Organizadores de workshops, escolas, simpósios e conferências;
\item Equipes que necessitam de uma agenda regional auditável e pesquisável.
\end{itemize}

\section{Princípios de projeto}

\begin{itemize}[leftmargin=*,itemsep=0.25em]
\item Fonte oficial antes de agregadores secundários;
\item Separação entre estado do evento, submissão e inscrição;
\item Preservação de divergências em vez de sobrescrita silenciosa;
\item Coleta respeitosa, com robots.txt, atraso e limitação de concorrência;
\item Rastreabilidade por URL, trecho de evidência, data de verificação e confiança;
\item Estrutura modular baseada no Princípio da Responsabilidade Única.
\end{itemize}
